\section{Discussion}
\label{sec:discussion}

\subsection{Why Did Dependency Milestone Scheduling Underperform?}
The central finding of this investigation—that fixed-periodic verification outperforms dependency-based milestone verification under matched budgets—contradicts intuitive assumptions prevalent in recent agent verification literature~\citep{afoud2026vreact}. To understand why milestone verification failed to contain error propagation, we must examine the topological structure of computational DAGs.

\paragraph{The Structural Clumping Phenomenon:}
In any valid multi-step computational DAG (including Fork-Join, Spine-Join, and Hierarchical Reduction topologies), a join node cannot execute until \emph{all} of its predecessor branches have fully completed their respective execution sequences. In practice, this creates an inherent topological skew: multi-parent integration boundaries naturally concentrate in the **middle and late phases** of the workflow horizon.

Consequently, milestone-based schedulers suffer from an \textbf{``early-stage blindspot''}:
\begin{itemize}
    \item In Arm B (Fixed-Periodic), the first verification check is scheduled unconditionally at $\lceil H/3 \rceil$ (step 4 at $H=12$, step 8 at $H=24$, step 12 at $H=36$). This enforces a strict temporal barrier across all active branches early in the workflow lifecycle.
    \item In Arm C (Milestone), the scheduler defers its first check until all inputs to Join 1 complete. In an $H=24$ graph, if Branch A requires 6 steps and Branch B requires 10 steps, Join 1 is not scheduled until step 10. Any corruption introduced at step 1 in Branch A propagates unmonitored through 9 subsequent operations before the join milestone is evaluated.
\end{itemize}

This topological clumping explains why the performance disparity between Arms B and C peaks during middle-stage faults ($\Delta \text{TCF} = +0.0960$). Under milestone scheduling, errors introduced in long parallel branches cascade through multiple intermediate transformations, contaminating peer branches and inflating both detection delay ($+2.31$ steps) and downstream taint consumption ($+3.91\%$).

\subsection{Re-Evaluating Architectural Intuition in Agent Workflows}
Our findings challenge the uncritical assumption that verification placement should be guided exclusively by semantic integration boundaries~\citep{afoud2026vreact,shoham2026jevact}. While dependency joins represent logical integration points where data streams compose, \textbf{logical integration points are not necessarily optimal risk-mitigation points}. 

When verification call budgets are severely constrained ($K \ll H$), uniform temporal coverage is fundamentally more resilient than topological placement. Fixed-periodic scheduling acts as a geometric bounding box, guaranteeing that an undetected error cannot propagate further than $\lceil H/K \rceil$ steps regardless of graph structure, branch asymmetry, or execution path length. In contrast, milestone scheduling couples check placement to graph geometry, rendering the agent vulnerable to long, unmonitored branch paths.

\subsection{Decoupling Containment from Task Completion}
A crucial insight from our empirical analysis is the complete decoupling between intermediate error containment (TCF) and final task completion ($38.46\%$ across all arms). 

Why did drastic improvements in detection latency ($5.50$ steps in Arm B vs. $12.17$ steps in Arm A) and taint containment ($0.8614$ in Arm B vs. $1.0000$ in Arm A) fail to increase final task completion?
\begin{enumerate}
    \item \textbf{Terminal Verification Parity:} Because Arm A evaluates the terminal artifact against the full contract suite at step $H$, any structurally detectable fault that reaches the end of the workflow is eventually identified.
    \item \textbf{Single-Replay Recovery Dynamics:} Under our controlled recovery protocol, whenever a fault is detected (whether at step 4 or step 12), the executor restores the last passing checkpoint and disarms the injector during replay. Because the fault is removed during re-execution, 100\% of detected episodes successfully recover, while 100\% of undetected episodes fail.
\end{enumerate}

Therefore, the primary architectural benefit of intermediate verification is not that it converts failed workflows into successful ones; rather, its true value lies in **containing blast radius and minimizing re-execution waste**. Arm B saved an average of $2.87$ re-executed nodes per episode relative to Arm A, cutting redundant compute cycles while shielding downstream external environments (e.g., databases, tool APIs, or multi-agent messaging channels) from corrupted state consumption.

\subsection{Architectural Implications for Compound AI Systems}
For software architects designing autonomous agent frameworks (e.g., LangGraph, AutoGen, ScienceFlow, and CrewAI), these results yield concrete guidelines:
\begin{enumerate}
    \item \textbf{Prioritize Periodic Cadences Under Low Budgets:} When verification budgets are limited ($K \le 3$), allocate checks at uniform periodic intervals rather than waiting for complex multi-agent joins.
    \item \textbf{Adopt Hybrid Two-Tier Scheduling:} If computational budgets permit $K > 3$, implement a hybrid policy: establish a fixed-periodic backbone to bound worst-case branch latency, supplemented by targeted milestone assertions immediately prior to state-merging operations.
    \item \textbf{Account for Replay Suffix Cost:} In workflows with high execution costs (e.g., wet-lab experiments or GPU-intensive training), checkpoint frequency should balance forward taint containment against the computational cost of suffix re-execution.
\end{enumerate}
